\documentclass[11pt]{article}

\usepackage[letterpaper,margin=1in]{geometry}
\usepackage[T1]{fontenc}
\usepackage{lmodern}
\usepackage{amsmath}
\usepackage{amssymb}
\usepackage{amsthm}
\usepackage{enumitem}
\usepackage{booktabs}
\usepackage[hidelinks]{hyperref}

\theoremstyle{remark}
\newtheorem{remark}{Remark}

\setlist[enumerate]{itemsep=1pt,topsep=3pt}

\newcommand{\cN}{\mathcal N}
\newcommand{\cK}{\mathcal K}
\newcommand{\cG}{\mathcal G}
\newcommand{\cP}{\mathcal P}
\newcommand{\cA}{\mathcal A}
\newcommand{\IO}{\mathrm{IO}}
\newcommand{\boxnote}[2]{\par\medskip\noindent\fbox{\parbox{\dimexpr\linewidth-2\fboxsep-2\fboxrule}{\textbf{#1.} #2}}\par\medskip}

\title{Solving the Model, Step by Step}
\author{}
\date{}

\begin{document}
\maketitle

\noindent This document derives the solution of the model in the companion model
section, one substitution at a time. It adds no new assumptions and no new results.
Its purpose is to show exactly how the equilibrium conditions are reduced to a small
system that can be solved, which objects have closed forms, and which ones must be
computed. Propositions and theorems cited by name and number are those of the model
section.

The derivation is written for the baseline: quantity competition, an upper-tier
elasticity $\eta=1$, intermediate inputs with shares $\nu_k\in[0,1)$, and non-negative
tariffs. Boxed notes say what changes under price competition and when $\eta>1$, which
the model section allows only without intermediate inputs.

\tableofcontents

\setcounter{section}{-1}
\section{What Has to Be Solved}
\label{sec:overview}

\paragraph{Primitives.} The economy is described by: the sets of countries $\cN$ and
sectors $\cK$; labor endowments $L_n$; preference weights $\beta_k$ and elasticities
$\sigma_k>\eta=1$; intermediate shares $\nu_k$ and input--output weights
$\omega_{k'k}$; plant productivities $\varphi_i$ (which already contain the location,
plant and group components and the head-office capability terms); the efficiency
losses $\gamma_{h\ell}$; trade costs $d_{\ell n}$; tariffs $t_{\ell nk}$; fixed-cost
parameters $f_k$ and $f^d_{\ell n}$; the pool of potential groups with their potential
plants $\cP_g$; and ownership shares $\theta_{gn}$. All of these are fixed before any
decision is taken.

\paragraph{Unknowns.} An equilibrium (Definition 2 of the model section) consists of
wages $w_n$; sectoral price indices $P_{nk}$; market sizes $E_{nk}$; incomes $X_n$;
active sets $T_{g,n}\subseteq\cP_g$, one per group and destination of its sector; and
the price $p_{i,n}$ and quantity $q_{i,n}$ of every active plant in every market it
serves.

\paragraph{Strategy.} There is no closed form for the equilibrium as a whole. What
the model allows is a chain of substitutions in which every object becomes a function
of two things only: the wage vector $w$ and the configuration of active sets
$T=\{T_{g,n}\}$. Each link of the chain is either a closed form or a problem with a
unique solution that is easy to compute. The equilibrium is then a pair $(w,T)$ that
satisfies $N-1$ labor-market conditions and an entry condition in every market. The
steps are:

\begin{center}
\small
\begin{tabular}{@{}llll@{}}
\toprule
Step & Solves for & Takes as given & How \\
\midrule
1 & demand functions & prices, spending & closed form \\
2 & unit and delivered costs & wages, input prices & closed form \\
3 & shares, markups, prices in one market & delivered costs, active plants & one scalar equation \\
4 & active plants in one market & costs, fixed costs, spending & ladder and bisection \\
5 & input prices & wages, active sets & contraction \\
6 & market sizes and incomes & wages, active sets, prices & linear system \\
7 & labor-market conditions & everything above & closed form \\
8 & wages & active sets & $N-1$ equations \\
9 & active sets & everything above & outer loop \\
\bottomrule
\end{tabular}
\end{center}

\noindent Steps 1--3 and 7 hold at any wages. Steps 3--6 are nested: market shares do
not depend on spending, so input prices can be found before market sizes, and market
sizes and incomes solve a linear system once prices are known.

\section{Step 1: Household Demand}
\label{sec:demand}

\paragraph{Lower tier.} Fix country $n$ and sector $k$, and write $\sigma=\sigma_k$.
For a given amount of the sectoral composite $C_{nk}$, the household minimizes
spending $\sum_{i\in\cA_{nk}}p_{i,n}q_{i,n}$ subject to
$\bigl[\sum_iq_{i,n}^{(\sigma-1)/\sigma}\bigr]^{\sigma/(\sigma-1)}=C_{nk}$. With
multiplier $\chi$, the first-order condition for variety $i$ is
\[
  p_{i,n}=\chi\,\frac{\partial C_{nk}}{\partial q_{i,n}}
  =\chi\,C_{nk}^{1/\sigma}q_{i,n}^{-1/\sigma},
  \qquad\text{so}\qquad
  q_{i,n}=\Bigl(\frac{p_{i,n}}{\chi}\Bigr)^{-\sigma}C_{nk}.
\]
Substituting into the constraint,
$C_{nk}=\chi^{\sigma}C_{nk}\bigl[\sum_ip_{i,n}^{1-\sigma}\bigr]^{\sigma/(\sigma-1)}$,
so $\chi=\bigl[\sum_ip_{i,n}^{1-\sigma}\bigr]^{1/(1-\sigma)}\equiv P_{nk}$. The
multiplier is the price index, and minimized spending is
$\sum_ip_{i,n}q_{i,n}=P_{nk}C_{nk}$. Hence
\begin{equation}
  q_{i,n}=\Bigl(\frac{p_{i,n}}{P_{nk}}\Bigr)^{-\sigma}C_{nk}
  =p_{i,n}^{-\sigma}P_{nk}^{\sigma-1}\,(P_{nk}C_{nk}).
  \label{eq:lowertier}
\end{equation}

\paragraph{Upper tier.} With $\eta=1$ the upper tier is Cobb--Douglas,
$U_n=\prod_kC_{nk}^{\beta_k}$, and maximizing it subject to
$\sum_kP_{nk}C_{nk}=X_n$ gives the expenditure shares $P_{nk}C_{nk}=\beta_kX_n$.
Final spending on sector $k$ in $n$ is therefore a fixed share of income,
\begin{equation}
  E^F_{nk}=\beta_kX_n .
  \label{eq:finalspending}
\end{equation}

\paragraph{Total demand.} Firms buy intermediate inputs as the same sectoral
composites (Step 2), so they face the same price index and split their purchases
across varieties by the same rule \eqref{eq:lowertier}. Adding households and firms,
and writing $E_{nk}$ for total spending on sector $k$ in $n$, demand for plant $i$'s
variety in $n$ is
\begin{equation}
  q_{i,n}=p_{i,n}^{-\sigma}P_{nk}^{\sigma-1}E_{nk},
  \qquad
  P_{nk}=\Bigl[\sum_{i\in\cA_{nk}}p_{i,n}^{1-\sigma}\Bigr]^{1/(1-\sigma)} .
  \label{eq:demand}
\end{equation}
The revenue share of plant $i$ in the market is
$s_{i,n}\equiv p_{i,n}q_{i,n}/E_{nk}=p_{i,n}^{1-\sigma}/P_{nk}^{1-\sigma}$, and these
shares add up to one.

\boxnote{When $\eta>1$}{Only without intermediates. The upper tier gives
$P_{nk}C_{nk}=\beta_k(P_{nk}/P_n)^{1-\eta}X_n$ with
$P_n=[\sum_k\beta_kP_{nk}^{1-\eta}]^{1/(1-\eta)}$, so final spending on a sector
rises when its relative price falls. Everything below goes through with $\beta_k$
replaced by $\beta_k(P_{nk}/P_n)^{1-\eta}$, which is known once prices are known
(Step 5), so Step 6 is still a linear system in incomes.}

\section{Step 2: Costs}
\label{sec:costs}

\paragraph{Unit cost.} Plant $i$ in sector $k$, located in $\ell$, combines labor
$l$ with a bundle $B$ of intermediate inputs, which in turn combines quantities
$m_{k'}$ of the sectoral composites bought in $\ell$:
\[
  y=\frac{\varphi_i}{\gamma_{h_i\ell}}
  \Bigl(\frac{l}{1-\nu_k}\Bigr)^{1-\nu_k}\Bigl(\frac{B}{\nu_k}\Bigr)^{\nu_k},
  \qquad
  B=\prod_{k'}\Bigl(\frac{m_{k'}}{\omega_{k'k}}\Bigr)^{\omega_{k'k}} .
\]
The constants in the denominators are the normalization under which the Cobb--Douglas
cost function has no extra constant; the model section absorbs this constant into
$\varphi_i$, which is equivalent. Minimizing $\sum_{k'}P_{\ell k'}m_{k'}$ for one unit
of $B$ gives the cost of the bundle, $P^{\IO}_{\ell k}=\prod_{k'}P_{\ell k'}^{\omega_{k'k}}$,
with spending on sector $k'$ equal to the share $\omega_{k'k}$ of bundle spending.
Minimizing $w_\ell l+P^{\IO}_{\ell k}B$ for one unit of output then gives the unit
cost
\begin{equation}
  a_i=\frac{\gamma_{h_i\ell}}{\varphi_i}\,w_\ell^{\,1-\nu_k}\bigl(P^{\IO}_{\ell k}\bigr)^{\nu_k},
  \qquad
  P^{\IO}_{\ell k}=\prod_{k'\in\cK}P_{\ell k'}^{\,\omega_{k'k}},
  \label{eq:unitcost}
\end{equation}
which is the cost equation of the model section. By Shephard's lemma, for every
dollar of production cost the plant pays $1-\nu_k$ to labor in $\ell$ and
$\nu_k\omega_{k'k}$ for the composite of sector $k'$ in $\ell$.

\paragraph{Delivered cost.} To sell $q_{i,n}$ units in $n$ the plant produces
$d_{\ell n}q_{i,n}$ units. Its variable cost of serving $n$, exclusive of tariffs, is
\begin{equation}
  V_{i,n}=a_id_{\ell n}q_{i,n},
  \label{eq:V}
\end{equation}
country $n$ collects $t_{\ell nk}V_{i,n}$, and the delivered marginal cost is
\begin{equation}
  c_{i,n}=a_i\,d_{\ell n}\,(1+t_{\ell nk}).
  \label{eq:delivered}
\end{equation}

\paragraph{Where the money goes.} Out of the variable cost $V_{i,n}$, labor in $\ell$
receives $(1-\nu_k)V_{i,n}$ and the market $(\ell,k')$ receives
$\nu_k\omega_{k'k}V_{i,n}$ as intermediate purchases. If the plant serves $n$, it
also pays the fixed cost
\begin{equation}
  F_{i,n}=f_k\,w_\ell\,f^d_{\ell n}
  \label{eq:fixed}
\end{equation}
to labor in $\ell$. These are the only payments a plant makes; what remains of its
revenue is tariff revenue and profit (Step 3).

\section{Step 3: One Market with Given Plants}
\label{sec:market}

Fix a market $(n,k)$ and the set of active plants $\cA$ in it, with $T_g\subseteq\cA$
the plants of group $g$. Drop the market's indices: $\sigma$, $E$, $P$, $c_i$, $s_i$,
$t_i\equiv t_{\ell(i)nk}$. The inputs of this step are the delivered costs $c_i$ and
the spending level $E$. Its outputs are prices, quantities, shares, markups and
profits.

\subsection{Inverse demand}

Let $\rho=(\sigma-1)/\sigma$ and $Q=(\sum_iq_i^\rho)^{1/\rho}$, the sectoral composite.
Since $PQ=E$, the lower-tier rule \eqref{eq:lowertier} can be inverted as
$p_i=P(q_i/Q)^{-1/\sigma}=(E/Q)(q_i/Q)^{-1/\sigma}$, so that
\begin{equation}
  p_i=E\,Q^{-\rho}q_i^{\rho-1},
  \qquad
  s_i=\frac{p_iq_i}{E}=\frac{q_i^\rho}{Q^\rho} .
  \label{eq:inverse}
\end{equation}
With $\eta=1$, $E$ is fixed from the point of view of the groups in the market.

\subsection{The group's first-order condition and the markup}

Group $g$ chooses $q_i$ for every $i\in T_g$ to maximize
$\pi_g=\sum_{i\in T_g}(p_i-c_i)q_i$. Its revenue is
\[
  R_g=\sum_{j\in T_g}p_jq_j=E\,Q^{-\rho}\sum_{j\in T_g}q_j^\rho=E\,S_g,
  \qquad
  S_g\equiv\sum_{j\in T_g}s_j .
\]
Using $\partial\ln Q/\partial q_i=q_i^{\rho-1}/Q^\rho=s_i/q_i$, the marginal revenue of
plant $i$ is
\[
  \frac{\partial R_g}{\partial q_i}
  =E\,Q^{-\rho}\Bigl[\rho\,q_i^{\rho-1}-\rho\,\frac{q_i^{\rho-1}}{Q^\rho}\sum_{j\in T_g}q_j^\rho\Bigr]
  =\rho\,p_i\,(1-S_g).
\]
The second term is cannibalization: expanding plant $i$ lowers the price index and
takes sales from every variety in the market, including the group's other plants, and
the group counts the loss on all of them, which is proportional to its total share
$S_g$. Setting marginal revenue equal to marginal cost $c_i$ gives
\begin{equation}
  p_i=\mu(S_g)\,c_i,
  \qquad
  \mu(S)=\frac{\sigma}{(\sigma-1)(1-S)} .
  \label{eq:markup}
\end{equation}
This is Proposition 1 of the model section at $\eta=1$. The markup is common to all
plants of the group and depends only on its total share. It rises from
$\mu(0)=\sigma/(\sigma-1)$ and is unbounded as $S\to1$, which is why every market needs
at least two active groups. The markup elasticity, used below, is
\begin{equation}
  \varepsilon_\mu(S)\equiv\frac{S\mu'(S)}{\mu(S)}=\frac{S}{1-S} .
  \label{eq:epsmu}
\end{equation}

\subsection{Shares within a group}

With a common markup, $p_i/p_j=c_i/c_j$ for any two plants of the same group, so by
\eqref{eq:demand} $s_i/s_j=(c_i/c_j)^{1-\sigma}$. Summing over the group,
\begin{equation}
  \frac{s_i}{S_g}=\frac{c_i^{1-\sigma}}{K_g},
  \qquad
  K_g\equiv\sum_{i\in T_g}c_i^{1-\sigma}.
  \label{eq:withingroup}
\end{equation}
Costs show up in how the group splits its sales across plants, not in markups.

\subsection{The share equation}

Let $A\equiv\sum_{i\in\cA}p_i^{1-\sigma}=P^{1-\sigma}$. Then $s_i=p_i^{1-\sigma}/A$.
Substituting $p_i=\mu(S_g)c_i$ and summing over the group,
\[
  S_g=\sum_{i\in T_g}\frac{\mu(S_g)^{1-\sigma}c_i^{1-\sigma}}{A}
  =\frac{\mu(S_g)^{1-\sigma}K_g}{A},
\]
that is,
\begin{equation}
  \psi(S_g)=\frac{K_g}{A},
  \qquad
  \psi(S)\equiv S\,\mu(S)^{\sigma-1}
  =\Bigl(\frac{\sigma}{\sigma-1}\Bigr)^{\sigma-1}\frac{S}{(1-S)^{\sigma-1}} .
  \label{eq:sharemap}
\end{equation}
For a given $A$ this is one equation in one unknown for each group. Taking logs and
moving the constant to the right,
\begin{equation}
  \ln S_g-(\sigma-1)\ln(1-S_g)=\ln K_g-\ln A-(\sigma-1)\ln\frac{\sigma}{\sigma-1}
  \equiv x_g(A).
  \label{eq:shareeq}
\end{equation}
The left side rises strictly from $-\infty$ to $+\infty$ as $S_g$ goes from $0$ to
$1$, so the equation has exactly one root $S_g(A)\in(0,1)$ for every $A>0$. It has no
closed form except in special cases, but it is a monotone scalar equation and Newton's
method converges quickly from any interior start, using the derivative
\[
  \frac{d}{dS}\bigl[\ln S-(\sigma-1)\ln(1-S)\bigr]
  =\frac{1+(\sigma-2)S}{S(1-S)} .
\]

\subsection{The market-clearing equation}

Shares must add up to one. With $S_g(A)$ the root of \eqref{eq:shareeq}, the
equilibrium of the market is the value of $A$ that solves
\begin{equation}
  \mathcal S(A)\equiv\sum_gS_g(A)=1 .
  \label{eq:clearing}
\end{equation}
Differentiating \eqref{eq:sharemap} in logs, $[1+(\sigma-1)\varepsilon_\mu(S_g)]\,
d\ln S_g=-d\ln A$, so
\begin{equation}
  \frac{dS_g}{d\ln A}=-S_g\Lambda_g,
  \qquad
  \Lambda_g\equiv\frac{1}{1+(\sigma-1)\varepsilon_\mu(S_g)}
  =\frac{1-S_g}{1+(\sigma-2)S_g}>0 .
  \label{eq:dSdA}
\end{equation}
Every share falls as $A$ rises, so $\mathcal S$ is strictly decreasing and
\eqref{eq:clearing} has at most one solution (Proposition 2). A bracket that contains
it is found as follows. Since $\mu(S)\ge\mu(0)$, \eqref{eq:sharemap} gives
$S_g\le K_g/[A\,\mu(0)^{\sigma-1}]$, so at
\[
  A_{\mathrm{hi}}=\mu(0)^{1-\sigma}\sum_gK_g
\]
shares add up to at most one. As $A\to0$ every root of \eqref{eq:shareeq} tends to
one, so $\mathcal S$ tends to the number of groups, which is at least two; halving $A$
from $A_{\mathrm{hi}}$ reaches a point $A_{\mathrm{lo}}$ with
$\mathcal S(A_{\mathrm{lo}})\ge1$ in finitely many steps. Between the two, Newton's
method on $\ln A$ with derivative $\mathcal S'(\ln A)=-\sum_gS_g\Lambda_g$, falling
back to bisection whenever a step would leave the bracket, converges to the unique
root.

\subsection{Recovering everything else}

Once $A$ and the shares $S_g$ are known, every other object in the market follows in
closed form:
\begin{align}
  \mu_g&=\mu(S_g),\qquad p_i=\mu_gc_i,\qquad P=A^{1/(1-\sigma)},
  \label{eq:recover1}\\
  s_i&=S_g\,\frac{c_i^{1-\sigma}}{K_g},\qquad r_i=s_iE,\qquad q_i=\frac{r_i}{p_i},
  \label{eq:recover2}\\
  V_i&=\frac{r_i}{\mu_g(1+t_i)},\qquad \text{tariff revenue } t_iV_i .
  \label{eq:recover3}
\end{align}
The price index is consistent with the prices, since
$\sum_ip_i^{1-\sigma}=\sum_g\mu_g^{1-\sigma}K_g=A\sum_gS_g=A$. Variable profit is
\[
  \pi_g=\sum_{i\in T_g}(p_i-c_i)q_i=\Bigl(1-\frac{1}{\mu_g}\Bigr)\sum_{i\in T_g}r_i
  =E\,S_g\,m(S_g),
  \qquad
  m(S)\equiv1-\frac{1}{\mu(S)}=\frac{1+(\sigma-1)S}{\sigma},
\]
so that
\begin{equation}
  \pi_g=\frac{E}{\sigma}\bigl(S_g+\kappa S_g^2\bigr),\qquad\kappa=\sigma-1 .
  \label{eq:profit}
\end{equation}
Profit does not depend on the tariff directly: the tariff enters through $c_i$, which
moves shares.

\paragraph{The revenue identity.} For each plant, revenue splits exactly into
variable cost, tariff revenue and profit:
\begin{equation}
  r_i=\mu_g(1+t_i)V_i=V_i+t_iV_i+(\mu_g-1)(1+t_i)V_i,
  \label{eq:revid}
\end{equation}
and $(\mu_g-1)(1+t_i)V_i=(p_i-c_i)q_i$ is plant $i$'s contribution to $\pi_g$. Per
dollar of spending in the market, the three parts are
\begin{equation}
  \underbrace{\sum_i\frac{s_i}{\mu_{g(i)}(1+t_i)}}_{\text{variable cost}}
  +\underbrace{\sum_i\frac{t_is_i}{\mu_{g(i)}(1+t_i)}}_{\text{tariffs}}
  +\underbrace{\sum_gS_g\,m(S_g)}_{\text{profit}}=1 .
  \label{eq:perdollar}
\end{equation}

\paragraph{Two properties used below.} First, shares, markups and the price index
depend only on the delivered costs, never on $E$: spending scales quantities, revenue,
costs and profits, nothing else. Second, multiplying every delivered cost by $\lambda$
multiplies every $K_g$ and $A$ by $\lambda^{1-\sigma}$, leaves every share and markup
unchanged, and multiplies every price and $P$ by $\lambda$. The price index is
therefore homogeneous of degree one in the delivered costs.

\boxnote{Under price competition}{Demand gives $\partial\ln q_i/\partial\ln p_j
=-\sigma\mathbf 1\{i=j\}+(\sigma-\eta)s_j$. Imposing a common markup in the group's
first-order conditions gives $\mu(S)=\varepsilon(S)/[\varepsilon(S)-1]$ with
$\varepsilon(S)=\sigma-(\sigma-\eta)S$ (Proposition 1). Steps 3.3--3.6 are unchanged
with this $\mu$: the within-group rule \eqref{eq:withingroup}, the share equation
$\psi(S_g)=K_g/A$ with $\psi(S)=S\mu(S)^{\sigma-1}$, the clearing equation and the
bracket all carry over, and $\Lambda_g=1/[1+(\sigma-1)\varepsilon_\mu(S_g)]$ with the
new $\varepsilon_\mu$. The Lerner index is $1/\varepsilon(S)$, so
$\pi_g=ES_g/\varepsilon(S_g)$, which is not quadratic in $S_g$.}

\boxnote{When $\eta>1$ (quantity competition)}{The markup is
$1/\mu(S)=1-(1-S)/\sigma-S/\eta$, which is bounded above by $\eta/(\eta-1)$, and
$\kappa=\sigma/\eta-1$. Because $\psi(1)=\mu(1)^{\sigma-1}$ is finite, the share
equation has an interior root only when $K_g/A<\psi(1)$; the clearing equation is
solved on $A>\max_gK_g/\psi(1)$, where $\mathcal S$ is still strictly decreasing.}

\section{Step 4: Entry into One Market}
\label{sec:entry}

Now let each group choose which of its potential plants serve the market. The inputs
are the delivered costs $c_i$ and fixed costs $F_i$ of all potential plants $i\in\cP_g$
and the spending level $E$. The outputs are the active sets $T_g$ and the aggregate
$A$.

\subsection{The value of a plant set}

By Step 3, a set $T\subseteq\cP_g$ affects the group's profit only through its strength
$K(T)=\sum_{i\in T}c_i^{1-\sigma}$, and costs it $\Phi(T)=\sum_{i\in T}F_i$. Taking
$A$ as given, the group's share with set $T$ is $\psi^{-1}(K(T)/A)$, and its net payoff
is
\begin{equation}
  U_g(T;A)=\Omega\bigl(K(T);A\bigr)-\Phi(T),
  \qquad
  \Omega(K;A)=\frac{E}{\sigma}\bigl(S+(\sigma-1)S^2\bigr)\Big|_{S=\psi^{-1}(K/A)} .
  \label{eq:payoff}
\end{equation}
The empty set gives $U_g=0$.

\subsection{The ladder}

A group with $|\cP_g|$ potential plants has $2^{|\cP_g|}$ sets. Since $\Omega$ is
strictly increasing in $K$, a set can be chosen only if no other set has at least as
much strength at no greater fixed cost. The construction is:
\begin{enumerate}
\item compute $(K(T),\Phi(T))$ for every $T\subseteq\cP_g$;
\item discard every $T$ for which some $T'$ has $K(T')\ge K(T)$ and
$\Phi(T')\le\Phi(T)$, with one inequality strict;
\item order the remaining sets by strength. Fixed cost rises along the order too,
otherwise a set would have been discarded. Call them rungs $0,1,\dots,R_g$, with rung
0 the empty set.
\end{enumerate}
When all of a group's plants pay the same fixed cost $F$ in the market, the ladder is
explicit: rung $r$ is the set of the $r$ plants with the lowest delivered costs, since
among sets with $r$ plants, which all cost $rF$, these have the most strength. A
stand-alone firm has a single plant and a ladder with two rungs, out or in.

\subsection{The best response at a given aggregate}

At a given $A$, the group's best response is the rung with the highest payoff,
\begin{equation}
  T^*_g(A)=\arg\max_{r=0,\dots,R_g}\;\Omega(K_r;A)-\Phi_r,
  \qquad
  K^*_g(A)=K\bigl(T^*_g(A)\bigr).
  \label{eq:bestresponse}
\end{equation}
In general this is a comparison of $R_g+1$ numbers. In the equal-fixed-cost case it
reduces to a cutoff. The strength added by each rung,
$c_{(r+1)}^{1-\sigma}$, falls with $r$ because plants are added from cheapest to most
expensive, and $\Omega$ is strictly concave in $K$ (Theorem 1(a)). So the gain from
climbing one more rung, $\Omega(K_{r+1};A)-\Omega(K_r;A)-F$, falls with $r$, and the
group climbs while this gain is positive.

\subsection{The entry equilibrium of the market}

An entry equilibrium (Definition 1) is a pair $(\{T_g\},A)$ in which every $T_g$ is a
best response at $A$ and $A$ clears the market with the chosen strengths. Substituting
the best responses into \eqref{eq:clearing} reduces it to one equation in $A$:
\begin{equation}
  \mathcal S^*(A)\equiv\sum_g\psi^{-1}\!\Bigl(\frac{K^*_g(A)}{A}\Bigr)=1 .
  \label{eq:entryclearing}
\end{equation}
When no group would hold more than half of the market with any of its rungs at the
aggregates compared, each $K^*_g(A)$ is non-increasing in $A$ (Corollary 1): a more
competitive market makes each group climb fewer rungs. Each term of $\mathcal S^*$
then falls with $A$ for two reasons, a lower strength and a higher aggregate, so
$\mathcal S^*$ is decreasing and crosses one at most once. It is solved by bisection
on $\ln A$, with the bracket of Step 3 computed from the strengths of the full sets
$\cP_g$ at the top and halving until $\mathcal S^*\ge1$ at the bottom, which is
possible when at least two groups can cover their fixed costs, the case the model
section considers.

\paragraph{The integer problem.} Because $K^*_g(A)$ moves in jumps, $\mathcal S^*$ can
jump over one. The bisection then converges to the aggregate $\bar A$ at which the
jump occurs, with $\mathcal S^*$ above one just below $\bar A$ and below one just above
it. At $\bar A$ some group switches rungs. If it is assigned its larger rung, the
aggregate that clears the market lies at or above $\bar A$, where the smaller rung is
its best response; if it is assigned its smaller rung, the clearing aggregate lies
below $\bar A$, where the larger rung is. Either way the group wishes to deviate, and no
exact equilibrium in whole plants exists. The procedure keeps the configuration with the fewest groups wishing to
deviate and reports how many there are.

\paragraph{Scale.} Multiplying every delivered cost, fixed cost and spending level by
the same $\lambda$ multiplies every payoff \eqref{eq:payoff} by $\lambda$ and leaves
shares unchanged, so it leaves every best response and the entry equilibrium
unchanged. In general equilibrium, this is why entry is unaffected by a common change
in all wages.

\section{Step 5: Input Prices, Given Wages and Active Sets}
\label{sec:prices}

From now on, wages $w$ and active sets $T$ are given. The delivered cost of every
plant depends on input prices through \eqref{eq:unitcost}, and input prices are the
price indices that those costs generate through Step 3. Substituting one into the
other gives a fixed point in the price indices.

\paragraph{The map.} Start from a vector of price indices $P=\{P_{\ell k}\}$ and
compute, in order:
\begin{enumerate}
\item the input bundle price $P^{\IO}_{\ell k}=\prod_{k'}P_{\ell k'}^{\omega_{k'k}}$ in
every country and sector;
\item the unit cost $a_i$ of every plant, by \eqref{eq:unitcost}, and its delivered
cost $c_{i,n}$ in every market it serves, by \eqref{eq:delivered};
\item in every market, the strengths $K_g$, the root $A$ of \eqref{eq:clearing} and the
new price index $A^{1/(1-\sigma)}$, by Step 3.
\end{enumerate}
Call the result $\mathcal T(P;w)$. Input prices are a fixed point,
$P=\mathcal T(P;w)$.

\paragraph{Why the fixed point is unique.} In logarithms, raising every price index by
$x\ge0$ raises $\ln P^{\IO}_{\ell k}$ by $x$, $\ln a_i$ by $\nu_kx$, and $\ln c_{i,n}$ by
$\nu_kx$. Since each market's price index is homogeneous of degree one and
non-decreasing in its delivered costs (Step 3), the new log price index of a
sector-$k$ market rises by at most $\nu_kx$. Hence
\[
  \ln\mathcal T(\ln P+x\mathbf 1)\le\ln\mathcal T(\ln P)+(\max_k\nu_k)\,x\mathbf 1,
\]
and the map is monotone. By Blackwell's sufficient conditions it is a contraction in
the sup norm with modulus $\max_k\nu_k<1$ (Proposition 3(a)). It has exactly one fixed
point, and iterating from any start converges to it: after $j$ iterations the error is
at most $(\max_k\nu_k)^j/(1-\max_k\nu_k)$ times the size of the first update.

\paragraph{Without intermediates.} When $\nu_k=0$ for all $k$, the unit cost is
$a_i=\gamma_{h_i\ell}w_\ell/\varphi_i$, delivered costs are known directly from wages,
and this step is a single pass through Step 3.

\paragraph{Output.} For given $(w,T)$, this step delivers every delivered cost
$c_{i,n}$, every share $s_{i,n}$ and $S_{g,n}$, every markup $\mu_{g,n}$ and every
price index $P_{nk}$. None of them depends on market sizes or incomes. Scaling all
wages by $\lambda$ scales every price index and every cost by $\lambda$ and leaves all
shares and markups unchanged.

\section{Step 6: Market Sizes and Incomes}
\label{sec:incomes}

With prices known, everything that remains is linear in spending. This step
substitutes Step 3 into the spending and income equations and solves the resulting
linear system.

\subsection{Spending as a function of incomes}

By \eqref{eq:recover3}, the variable cost of plant $i$ in market $(n',k')$ is a known
fraction of that market's spending:
\[
  V_{i,n'}=\frac{s_{i,n'}}{\mu_{g(i),n'}(1+t_{\ell(i)n'k'})}\,E_{n'k'} .
\]
Collect, for each country $n$, the part of market $(n',k')$'s spending that becomes
variable cost of plants located in $n$:
\begin{equation}
  v_{n,(n',k')}\equiv\sum_{i\in\cA_{n'k'}:\,\ell(i)=n}
  \frac{s_{i,n'}}{\mu_{g(i),n'}(1+t_{nn'k'})} .
  \label{eq:vcoef}
\end{equation}
Plants of sector $k'$ located in $n$ spend the share $\nu_{k'}\omega_{kk'}$ of their
variable costs on sector $k$ in $n$ (Step 2). Spending on market $(n,k)$ is final
spending \eqref{eq:finalspending} plus these purchases:
\begin{equation}
  E_{nk}=\beta_kX_n+\sum_{k'\in\cK}\omega_{kk'}\nu_{k'}
  \sum_{n'\in\cN}v_{n,(n',k')}\,E_{n'k'} .
  \label{eq:spending}
\end{equation}
Stack the $NK$ market sizes in a vector $E$ and the $N$ incomes in $X$. Define the
$NK\times N$ matrix $\mathcal E$ with entries $\mathcal E_{(n,k),j}=\beta_k\mathbf
1\{j=n\}$ and the $NK\times NK$ matrix $B$ with entries
$B_{(n,k),(n',k')}=\omega_{kk'}\nu_{k'}v_{n,(n',k')}$. Then
\eqref{eq:spending} reads
\begin{equation}
  E=\mathcal EX+BE
  \quad\Longrightarrow\quad
  E=(I-B)^{-1}\mathcal E\,X\equiv G\,X .
  \label{eq:EofX}
\end{equation}
The inverse exists and is non-negative. The column of $B$ for market $(n',k')$ sums to
\[
  \sum_{n,k}\omega_{kk'}\nu_{k'}v_{n,(n',k')}
  =\nu_{k'}\sum_{i\in\cA_{n'k'}}\frac{s_{i,n'}}{\mu_{g(i),n'}(1+t_{\ell(i)n'k'})}
  <\nu_{k'}<1,
\]
using $\sum_k\omega_{kk'}=1$ and the fact that markups exceed one and tariffs are
non-negative. A non-negative matrix whose columns sum to less than one has spectral
radius below one, so $(I-B)^{-1}=\sum_{s\ge0}B^s\ge0$. The entry $G_{(n,k),j}$ is the
spending on market $(n,k)$ generated, directly and through all rounds of intermediate
purchases, by one dollar of final spending by country $j$.

\subsection{Incomes as a function of spending}

Income is labor income plus owned net profits plus tariff revenue:
\[
  X_n=w_nL_n+\sum_g\theta_{gn}\Pi_g+T_n .
\]
Write each term in terms of spending. Group $g$ of sector $k$ earns variable profit
$E_{n'k}S_{g,n'}m(S_{g,n'})$ in each destination $n'$ by \eqref{eq:profit}, and pays
the fixed costs of its active plants,
\[
  \Phi_g\equiv\sum_{n'}\sum_{i\in T_{g,n'}}f_k\,w_{\ell(i)}\,f^d_{\ell(i)n'},
  \qquad
  \Pi_g=\sum_{n'}E_{n'k}S_{g,n'}m(S_{g,n'})-\Phi_g .
\]
The fixed costs depend on wages and active sets only. Tariff revenue in $n$ is, by
\eqref{eq:recover3},
\[
  T_n=\sum_k\tau_{nk}E_{nk},
  \qquad
  \tau_{nk}\equiv\sum_{i\in\cA_{nk}}\frac{t_{\ell(i)nk}\,s_{i,n}}{\mu_{g(i),n}(1+t_{\ell(i)nk})} .
\]
Define the $N\times NK$ matrix $H$ with entries
\begin{equation}
  H_{n,(n',k)}=\sum_{g:\,k(g)=k}\theta_{gn}S_{g,n'}m(S_{g,n'})+\mathbf 1\{n'=n\}\,\tau_{nk},
  \label{eq:H}
\end{equation}
the share of market $(n',k)$'s spending that becomes income of country $n$ as profit
or tariff revenue. Then
\[
  X=wL-\theta'\Phi+H\,E .
\]

\subsection{Solving the system}

Substituting \eqref{eq:EofX},
\begin{equation}
  X=wL-\theta'\Phi+\mathcal M X,
  \qquad
  \mathcal M\equiv HG,
  \qquad\Longrightarrow\qquad
  X=(I-\mathcal M)^{-1}\bigl(wL-\theta'\Phi\bigr),
  \label{eq:XofW}
\end{equation}
and then $E=GX$. Here $wL$ is the vector with entries $w_nL_n$ and $\theta'\Phi$ the
vector with entries $\sum_g\theta_{gn}\Phi_g$.

The inverse exists because every column of $\mathcal M$ sums to less than one.
Consider one dollar of final spending by country $j$. It generates the spending $Ge_j$
across markets. By \eqref{eq:perdollar}, the part of each market's spending that is
neither profit nor tariff revenue is variable cost. So the profits and tariffs
generated, which is the column sum $\mathbf 1'\mathcal Me_j$, equal total spending
generated minus total variable cost generated. Total spending generated is the one
dollar of final spending plus all intermediate purchases, and intermediate purchases
are the share $\nu$ of variable costs. Hence
\[
  \mathbf 1'\mathcal Me_j=1-\sum_{i}(1-\nu_{k(i)})\,\bigl[\text{variable cost of plant }i
  \text{ generated by the dollar}\bigr]<1,
\]
since the generated variable wage bill is positive when $\nu_k<1$. With $\mathcal M\ge0$
and column sums below one, $\rho(\mathcal M)<1$ and the solution is unique
(Proposition 3(b)).

\paragraph{Without intermediates.} With $\nu_k=0$, $B=0$, $G=\mathcal E$ and
$E_{nk}=\beta_kX_n$. The income system becomes $N$ equations,
\[
  X_n=w_nL_n-\sum_g\theta_{gn}\Phi_g
  +\sum_{n'}\Bigl[\sum_g\theta_{gn}\beta_{k(g)}S_{g,n'}m(S_{g,n'})\Bigr]X_{n'}
  +\Bigl[\sum_k\beta_k\tau_{nk}\Bigr]X_n .
\]

\paragraph{Output.} For given $(w,T)$, Steps 5 and 6 deliver every market size
$E_{nk}$, every income $X_n$, every group's net profit $\Pi_g$ and every country's
tariff revenue $T_n$, and through \eqref{eq:recover2}--\eqref{eq:recover3} every
quantity and variable cost. Scaling all wages by $\lambda$ scales the right side of
\eqref{eq:XofW} by $\lambda$ and leaves $\mathcal M$ unchanged, so incomes and market
sizes scale by $\lambda$ too.

\section{Step 7: The Labor Market}
\label{sec:labor}

\paragraph{Labor demand.} Country $n$'s wage bill is what plants located in $n$ pay
to labor: the share $1-\nu_k$ of their variable costs, plus the fixed costs of the
markets they serve. Using \eqref{eq:vcoef},
\begin{equation}
  W_n(w;T)=\sum_{k}(1-\nu_k)\sum_{n'}v_{n,(n',k)}E_{n'k}
  +\sum_{n'}\sum_{i:\,\ell(i)=n,\ i\in T_{g(i),n'}}f_{k(i)}\,w_n\,f^d_{nn'} .
  \label{eq:wagebill}
\end{equation}
Every term is known from Steps 5 and 6. The excess demand for labor is
\begin{equation}
  z_n(w;T)=\frac{W_n(w;T)}{w_n}-L_n,
  \label{eq:excess}
\end{equation}
and labor markets clear when $z_n=0$ for every $n$.

\paragraph{Walras' law.} One of these conditions is redundant: at any wages,
\[
  \sum_nw_nz_n(w;T)=\sum_nW_n-\sum_nw_nL_n=0 .
\]
To see it, add the revenue identity \eqref{eq:revid} over all plants and markets:
\[
  \sum_{\text{markets}}E_{nk}=\sum V+\sum_nT_n+\sum_g\pi_g ,
\]
where $\sum V$ is total variable cost and $\pi_g$ is the group's total variable profit.
Total spending is final spending, $\sum_nX_n$, plus intermediate purchases,
$\sum\nu V$. So $\sum_nX_n=\sum(1-\nu)V+\sum_nT_n+\sum_g\pi_g$. Adding the income
equations over countries, with $\sum_n\theta_{gn}=1$,
$\sum_nX_n=\sum_nw_nL_n+\sum_g(\pi_g-\Phi_g)+\sum_nT_n$. Equating the two expressions,
\[
  \sum_nw_nL_n=\sum(1-\nu)V+\sum_g\Phi_g=\sum_nW_n,
\]
since every fixed cost is paid to labor in the plant's location.

\paragraph{The numeraire.} Scaling all wages by $\lambda$ scales costs, prices and fixed
costs by $\lambda$ (Step 5), incomes and market sizes by $\lambda$ (Step 6), and leaves
shares, markups and active sets unchanged (Step 4). By \eqref{eq:wagebill}, $W_n$ scales
by $\lambda$, so $z$ is homogeneous of degree zero. Only relative wages are determined,
and we set $w_1=1$. By Walras' law, $z_1=0$ whenever $z_2=\dots=z_N=0$.

\section{Step 8: Wages, Given Active Sets}
\label{sec:wages}

\paragraph{The reduced system.} Steps 5--7 turn the whole model, for given active sets,
into $N-1$ equations in $N-1$ unknowns:
\begin{equation}
  z_n(1,w_2,\dots,w_N;T)=0,\qquad n=2,\dots,N .
  \label{eq:reduced}
\end{equation}
Evaluating $z(w;T)$ at a given $w$ means: solve the price fixed point (Step 5), build
$B$, $H$ and $\Phi$ and solve the linear system (Step 6), and compute the wage bills
\eqref{eq:wagebill}.

\paragraph{Fully written out without intermediates.} When $\nu_k=0$ the reduced system
can be written explicitly in terms of wages and market shares. Delivered costs are
$c_{i,n}=\gamma_{h_i\ell}w_\ell d_{\ell n}(1+t_{\ell nk})/\varphi_i$, so shares in every
market depend on wages only through relative delivered costs. Spending is
$E_{nk}=\beta_kX_n$, and incomes solve the $N$-equation system of Step 6. The wage
bill is
\[
  W_n=\sum_{n'}\sum_k\beta_kX_{n'}
  \sum_{i\in\cA_{n'k}:\,\ell(i)=n}\frac{s_{i,n'}}{\mu_{g(i),n'}(1+t_{nn'k})}
  +w_n\sum_{n'}\sum_{i:\,\ell(i)=n,\ i\in T_{g(i),n'}}f_{k(i)}f^d_{nn'} ,
\]
and \eqref{eq:reduced} requires $W_n=w_nL_n$. A rise in a foreign wage $w_j$ affects
$W_n$ in three ways: through the shares, as plants located in $n$ gain sales from
plants located in $j$; through the markups that the share changes imply; and through
incomes $X_{n'}$, which include the profits of the groups each country owns. The last
channel is the ownership channel of the model.

\paragraph{Damped tâtonnement.} The solution algorithm updates wages by
\begin{equation}
  w_n^{(j+1)}=\bigl(w_n^{(j)}\bigr)^{1-\lambda}\Bigl(\frac{W_n(w^{(j)};T)}{L_n}\Bigr)^{\lambda},
  \qquad\text{that is,}\qquad
  \ln w_n^{(j+1)}=\ln w_n^{(j)}+\lambda\ln\Bigl(1+\frac{z_n(w^{(j)};T)}{L_n}\Bigr),
  \label{eq:tatonnement}
\end{equation}
and then rescales so that $w_1=1$, which is allowed because the map is homogeneous of
degree one. A country's wage rises when its labor is in excess demand. A wage vector
is a fixed point exactly when $W_n=w_nL_n$ for every $n$, so the fixed points are the
equilibrium wages.

\paragraph{How much damping.} Let $M_{nj}=\partial\ln W_n/\partial\ln w_j$. In
logarithms the map \eqref{eq:tatonnement} has Jacobian $(1-\lambda)I+\lambda M$, whose
rows sum to one because $W$ is homogeneous of degree one. Its off-diagonal entries are
positive when labor is a gross substitute, $M_{nj}>0$ for $n\ne j$. Its diagonal
entries are positive when
\[
  \lambda<\frac{1}{1-\min_nM_{nn}} .
\]
When labor demand is elastic, $M_{nn}<0$, the bound is below one, and an undamped
update ($\lambda=1$) can overshoot. When both conditions hold on a compact set of wages
that is convex in logarithms, the map is a contraction in Hilbert's projective metric
on that set (Theorem 2), and the iteration converges to the unique equilibrium in it.

\paragraph{Newton steps.} Close to the solution, Newton's method on
$F_n(\ln w)=\ln W_n-\ln w_n-\ln L_n$, $n=2,\dots,N$, converges faster. Its Jacobian is
\[
  \frac{\partial F_n}{\partial\ln w_j}=M_{nj}-\mathbf 1\{n=j\},
\]
restricted to $n,j\ge2$. The matrix $M$ can be obtained column by column, by perturbing
one $\ln w_j$ and re-evaluating Steps 5--7. The step is
$\Delta\ln w=-[M-I]_{-1,-1}^{-1}F$. Iterations stop when
$\max_n|z_n|/L_n$ falls below a tolerance.

\section{Step 9: Active Sets}
\label{sec:activesets}

Steps 5--8 hold the active sets fixed. In equilibrium, the active sets in every market
must form an entry equilibrium (Step 4) at the delivered costs, fixed costs and market
sizes that the wages generate. These depend on the active sets in all markets: through
input prices, which affect costs everywhere, and through incomes, which affect spending
everywhere. Active sets are therefore found by an outer loop.

\begin{enumerate}
\item Start from a configuration $T^{(0)}$, for example every potential plant active
in every market of its sector.
\item Given $T^{(j)}$, solve Steps 5--8 for wages, and record in every market the
delivered costs $c_{i,n}$, the fixed costs $F_{i,n}$, the market size $E_{nk}$ and the
aggregate $A_{nk}$.
\item In every market, compute each group's ladder and its best response at the
recorded aggregate (Steps 4.2--4.3), and its gain
$\Delta_{g,n}=U_g(T^*_{g,n};A_{nk})-U_g(T^{(j)}_{g,n};A_{nk})\ge0$.
\item If no group gains, stop: every group's set is a best response at its market's
aggregate, and the aggregate clears the market with those sets, so every market is in
an entry equilibrium and $(w,T^{(j)})$ is an equilibrium.
\item Otherwise, let the groups with the largest gains switch to their best responses,
call the result $T^{(j+1)}$, and return to step 2.
\end{enumerate}

\noindent Moving only the groups with the largest gains is a way of damping the outer
loop. Moving every group at once can make the configuration oscillate, because wages
computed for a dense configuration differ substantially from wages for a sparse one.
The choice affects the path of the algorithm, not its fixed points, and it does not
impose an order of entry. The loop is not guaranteed to converge. When it does not
reach a configuration in which no group wishes to move, it returns the one with the
fewest such groups and reports how many there are. Because the joint equilibrium of
wages and active sets is not known to be unique, the whole procedure is run from
several starting points and the resulting configurations are compared.

\section{The Complete Algorithm, and What to Check}
\label{sec:summary}

Putting the steps together, from the innermost loop outward:

\begin{enumerate}[label=(\Alph*)]
\item \emph{One market, plants given (Step 3).} Given delivered costs: solve each
share equation \eqref{eq:shareeq} for $S_g(A)$, and the clearing equation
\eqref{eq:clearing} for $A$; recover $\mu_g$, $p_i$, $P$ and $s_i$ by
\eqref{eq:recover1}--\eqref{eq:recover2}.
\item \emph{Input prices (Step 5).} Given $(w,T)$: iterate $P\mapsto\mathcal T(P;w)$,
which calls (A) in every market, until it converges.
\item \emph{Spending and incomes (Step 6).} Given $(w,T)$ and the output of (B): build
$B$, $H$ and $\Phi$; compute $X$ by \eqref{eq:XofW} and $E=GX$.
\item \emph{Labor market (Step 7).} Compute $W_n$ by \eqref{eq:wagebill} and
$z_n$ by \eqref{eq:excess}.
\item \emph{Wages (Step 8).} Update $w$ by \eqref{eq:tatonnement}, then by Newton
steps, repeating (B)--(D), until $z(w;T)=0$.
\item \emph{Active sets (Step 9).} At the wages from (E), compute every group's best
response in every market; if some group gains, update the configuration and return to
(B).
\end{enumerate}

\begin{center}
\small
\begin{tabular}{@{}ll@{}}
\toprule
Closed form & Computed \\
\midrule
demand, unit and delivered costs & each group's share, given $A$ (monotone scalar equation) \\
markups, prices, quantities given shares & each market's $A$ (monotone scalar equation) \\
within-group shares, variable profits & each group's ladder and best response (finite comparison) \\
revenue, cost and tariff splits & input prices (contraction) \\
wage bills given spending and shares & spending and incomes (linear system) \\
 & wages ($N-1$ equations) and active sets (outer loop) \\
\bottomrule
\end{tabular}
\end{center}

\paragraph{Checks at the solution.} Each of the following holds exactly at an
equilibrium and is a useful test of any implementation.
\begin{enumerate}
\item In every market, shares add up to one and $\sum_ip_i^{1-\sigma}=P^{1-\sigma}$.
\item For every plant, revenue equals variable cost plus tariff plus profit,
\eqref{eq:revid}.
\item Every country's income equals labor income plus owned net profits plus tariff
revenue, and world income equals world final spending.
\item Every labor market clears, and Walras' law holds also away from the solution,
which tests the accounting at arbitrary wages.
\item No group gains from changing its set at its market's aggregate.
\item The hypotheses of the uniqueness results: no group holds more than half of any
market (Theorem 1), and the wage-bill elasticities satisfy the gross-substitutes and
damping conditions of Theorem 2.
\end{enumerate}

\end{document}
